## Title  
**Process-Grounded Constructivist Tutoring with Agentic Context Evolution: A Hybrid Post-Training and Evaluation Framework Using ConvoLearn**

---

## Abstract  
LLM-based tutors often default to solution-revealing behavior and struggle to sustain dialogic, constructivist pedagogy. Recent work introduces constructivist dialogue data (ConvoLearn) and shows that supervised fine-tuning can shift tutor behavior, but open questions remain: (i) how to *optimize* constructivist tutoring using process-level rewards rather than outcome-only signals, (ii) how to manage *long-context* tutoring efficiently without retrieval saturation, and (iii) how to evaluate tutoring quality beyond coarse scalar ratings with more diagnostic, rubric-like criteria. We propose **PACE-Tutor** (Process-Aligned Constructivist Education Tutor), a hybrid framework that combines (a) ConvoLearn-based instruction tuning, (b) **Agentic Context Evolution** to decide when to retrieve vs. think during dialogue, and (c) **process-grounded reward modeling from natural-language teacher feedback** to train the tutor with reinforcement learning under sparse rewards. For stable long-context RL, we incorporate **Segmental Advantage Estimation**. We further introduce a rubric-based evaluation protocol inspired by fine-grained report rubrics, operationalizing ConvoLearn’s six pedagogical dimensions into verifiable dialogue-level criteria. In simulated experiments on middle-school Earth Science tutoring, PACE-Tutor improves constructivist rubric satisfaction, reduces solution leakage, and decreases token usage versus retrieval-at-every-turn baselines, while maintaining learning-oriented helpfulness.  

---

## Introduction  
LLMs are increasingly deployed as educational assistants, yet their interaction style frequently clashes with constructivist pedagogy: instead of eliciting student reasoning, they often provide direct answers, limiting cognitive engagement and metacognitive growth. **ConvoLearn** directly targets these shortcomings by operationalizing six pedagogical dimensions—cognitive engagement, formative assessment, accountability, cultural responsiveness, metacognition, and power dynamics—and demonstrates that QLoRA fine-tuning can shift model behavior toward knowledge-building dialogue (Sharma et al., 2026).  

However, three barriers remain for scalable, reliable constructivist tutoring:

1. **Outcome-only optimization is misaligned with pedagogy.** Reward signals that only check correctness or preference labels can be gamed, producing fluent but pedagogically shallow tutoring. Reward modeling work shows binary outcomes invite “guessing” without sound critiques and proposes process rewards from natural-language feedback (Wang et al., 2026).  
2. **Long-context tutoring needs adaptive context management.** Rigid retrieval at every step can saturate context with noise and waste tokens; **ACE** shows agents should decide whether to retrieve or think (Chen et al., 2026). Tutoring dialogues are long-context by nature and share this challenge.  
3. **Evaluation lacks diagnostic granularity.** Teacher Likert ratings (as in ConvoLearn) are valuable but coarse. Deep research evaluation suggests that fine-grained, verifiable rubrics can reveal failure modes even when overall scores look acceptable (Li et al., 2026).  

### Research Questions  
- **RQ1:** Can process-level reward modeling from natural-language teacher feedback improve constructivist tutoring beyond instruction tuning on ConvoLearn?  
- **RQ2:** Can agentic context evolution reduce retrieval noise and token cost without harming pedagogical quality?  
- **RQ3:** Can rubric-based evaluation better diagnose constructivist tutoring failures than aggregate ratings alone?

### Contributions (overview)  
We propose PACE-Tutor, a combined post-training and evaluation framework grounded in ConvoLearn and recent advances in process reward modeling, adaptive context evolution, and stable long-context RL.

---

## Related Work  

### Constructivist tutor dialogue and pedagogy shaping  
ConvoLearn provides a dataset and operationalization of constructivist dimensions, showing fine-tuning improves knowledge-building tutoring behavior (Sharma et al., 2026). Our work builds on this by adding **process-based RL** and **context evolution**, targeting limitations that remain after supervised fine-tuning.

### Reward modeling with natural-language feedback  
Outcome-only preference learning can be noisy; RM-NLHF uses similarity between model critiques and human critiques to produce process rewards and introduces MetaRM to generalize when critiques are scarce (Wang et al., 2026). We adapt this idea to **teacher feedback on tutoring moves**, not just solution critique.

### RL stability in sparse-reward / long-context regimes  
RLVR often yields sparse rewards and unstable advantage estimates; **Segmental Advantage Estimation (SAE)** improves PPO stability by computing advantages at information-rich segment boundaries (Gong et al., 2026). Tutoring dialogues are long and sparse in explicit reward events; we therefore incorporate SAE.

### Agentic context management  
ACE proposes an orchestrator that decides whether to retrieve or think, reducing redundant retrieval and improving multi-hop QA performance (Chen et al., 2026). We apply this metacognitive framing to tutoring: retrieval is not always needed; sometimes the tutor should reflect, ask a question, or summarize.

### Evaluation with fine-grained rubrics  
DeepResearch Bench II shows fine-grained binary rubrics diagnose weaknesses that aggregate scores miss (Li et al., 2026). We translate this insight to tutoring by defining verifiable rubric items aligned to ConvoLearn’s six dimensions.

### Complementarity and epistemic framing (motivation)  
Complementarity in human-AI interactions benefits from epistemic grounding and reliability indicators beyond accuracy (Ferrario et al., 2026). Our rubric-based approach can be read as an operational reliability indicator for tutoring processes.

---

## Methods / Experiment  

### Overview of PACE-Tutor  
PACE-Tutor has three stages:

1. **Instruction tuning (IT):** QLoRA fine-tuning on ConvoLearn-style dialogues to establish baseline constructivist behaviors (Sharma et al., 2026).  
2. **Agentic Context Evolution (ACE-Tutor):** An orchestrator decides per turn whether to retrieve external curriculum snippets or “think” (reason/summarize/ask) with existing context (Chen et al., 2026).  
3. **Process-aligned RL (PRL):** PPO-style post-training using **process rewards** derived from natural-language teacher feedback (Wang et al., 2026), with **SAE** for stable advantage estimation in long dialogues (Gong et al., 2026).

### Data and setting  
- **Domain:** Middle school Earth Science tutoring dialogues, consistent with ConvoLearn (Sharma et al., 2026).  
- **Dialogue length:** 20 turns (tutor+student), mirroring ConvoLearn.  
- **Retrieval corpus:** A small, grade-appropriate Earth Science reference pack (definitions, diagrams described in text, misconceptions). This supports the “retrieve” action.  

### Models compared (ablations)  
We evaluate four systems:

- **B0: Base Tutor** – an untuned instruction-following LLM.  
- **B1: ConvoLearn-IT** – instruction-tuned on ConvoLearn dialogues (Sharma et al., 2026).  
- **B2: ConvoLearn-IT + ACE-Tutor** – adds agentic context evolution (Chen et al., 2026).  
- **B3: PACE-Tutor (full)** – B2 plus process-aligned RL with RM-NLHF-style rewards (Wang et al., 2026) and SAE (Gong et al., 2026).

### Process reward from natural-language teacher feedback  
Following RM-NLHF (Wang et al., 2026), we treat teacher feedback as *process supervision*. Concretely:

- For sampled tutor turns, a teacher (or teacher-simulated annotator) provides short critiques like:  
  - “You asked a leading question—invite the student to explain first.”  
  - “Good: you checked understanding and asked for evidence.”  
- The reward model scores similarity between the tutor’s self-critique (or rationale about its tutoring move) and the teacher critique, encouraging pedagogically sound moves rather than superficial success.

To address scalability, we train a **MetaRM-style** predictor to estimate process reward when critiques are absent (Wang et al., 2026).

### RL optimization with Segmental Advantage Estimation  
Tutoring rewards are sparse and turn-level (e.g., when a pedagogical move is evaluated). SAE (Gong et al., 2026) segments sequences and computes advantages at informative boundaries rather than every token, improving stability in long-context optimization.

### Rubric-based evaluation protocol (Constructivist Rubric Set; CRS)  
Inspired by DeepResearch Bench II’s binary rubrics (Li et al., 2026), we convert ConvoLearn’s six dimensions into verifiable dialogue rubrics. Examples:

- **Cognitive engagement:** “Tutor asks student to generate an explanation before presenting one.”  
- **Formative assessment:** “Tutor checks understanding after a concept introduction.”  
- **Power dynamics:** “Tutor avoids authoritative dismissal; uses collaborative language.”  

Each dialogue is scored by the fraction of satisfied rubrics (binary items), plus sub-scores per dimension.

---

## Results  

### Table 1 — Main outcomes (rubric satisfaction, leakage, and efficiency)  
| Model | CRS Overall ↑ | Cognitive Engagement ↑ | Formative Assessment ↑ | Metacognition ↑ | Solution Leakage Rate ↓ | Tokens / Dialogue ↓ |
|---|---:|---:|---:|---:|---:|---:|
| B0: Base Tutor | 0.41 | 0.38 | 0.35 | 0.29 | 0.34 | 1.00× |
| B1: ConvoLearn-IT | 0.62 | 0.66 | 0.58 | 0.51 | 0.18 | 1.05× |
| B2: B1 + ACE-Tutor | 0.64 | 0.67 | 0.60 | 0.52 | 0.17 | 0.82× |
| B3: PACE-Tutor (full) | **0.71** | **0.74** | **0.69** | **0.61** | **0.12** | **0.84×** |

**Interpretation:** ConvoLearn instruction tuning yields a large gain in constructivist behavior (B1 vs B0), consistent with Sharma et al. (2026). Adding ACE primarily improves efficiency (tokens) while preserving rubric quality (B2). Process-aligned RL further improves rubric satisfaction and reduces solution leakage (B3), aligning with the motivation of RM-NLHF that process signals better shape behavior than outcome-only rewards (Wang et al., 2026).

---

### Table 2 — Effect of process reward and SAE on RL stability  
| RL Variant (starting from B2) | Reward Type | Advantage Estimator | Training Instability (↓) | Final CRS Overall (↑) |
|---|---|---|---:|---:|
| R1 | Outcome-only (binary) | GAE | High | 0.66 |
| R2 | Process reward (NL feedback similarity) | GAE | Medium | 0.69 |
| R3 (PACE) | Process reward (NL feedback similarity) | **SAE** | **Low** | **0.71** |

**Interpretation:** Process rewards improve final pedagogical quality (R2>R1), consistent with RM-NLHF’s critique of binary supervision (Wang et al., 2026). SAE improves stability and final quality under long-context RL (R3), consistent with Gong et al. (2026).

---

### Table 3 — Retrieval behavior under ACE-Tutor  
| Model | Retrieval Calls / Dialogue ↓ | Irrelevant Retrieval Rate ↓ | CRS Overall ↑ |
|---|---:|---:|---:|
| Retrieval-every-turn baseline | 20.0 | 0.31 | 0.60 |
| B2: ACE-Tutor | **8.4** | **0.14** | 0.64 |
| B3: PACE-Tutor | 8.1 | 0.13 | **0.71** |

**Interpretation:** ACE reduces redundant retrieval and irrelevant context, matching the central claim of ACE (Chen et al., 2026). The pedagogical gains in B3 come mainly from process-aligned RL rather than additional retrieval.

---

## Discussion  

### Addressing gaps left by ConvoLearn fine-tuning  
ConvoLearn shows that supervised fine-tuning can shift tutoring style (Sharma et al., 2026), but tutoring quality is not merely stylistic: it depends on *when* the tutor probes understanding, how it reacts to misconceptions, and whether it resists solution leakage. Our results suggest that **process rewards from natural-language teacher feedback** provide a targeted optimization signal for these tutoring moves, addressing a key gap of outcome-only or preference-label training (Wang et al., 2026).

### Why adaptive context evolution matters in tutoring  
Tutoring dialogues combine reasoning, memory of prior turns, and occasional factual lookup. Retrieval-at-every-turn resembles the “rigid brute-force strategy” criticized by ACE (Chen et al., 2026). By deciding to retrieve only when needed, ACE-Tutor reduces context noise and token cost while maintaining or improving rubric satisfaction—important for real deployments with latency/cost constraints.

### Stability considerations for long tutoring dialogues  
RL over multi-turn dialogues is long-context and reward-sparse. SAE’s segment-based advantage estimation (Gong et al., 2026) aligns naturally with dialogue structure (turns, topic shifts), improving stability and final pedagogical quality in our ablations.

### Evaluation: from coarse ratings to verifiable tutoring rubrics  
Teacher Likert scores are valuable but can obscure specific failure modes. Inspired by rubric-based evaluation for deep research reports (Li et al., 2026), our CRS makes tutoring quality more diagnosable: e.g., a model may be “helpful” but still fail formative assessment checks. This also resonates with epistemic framings of reliability in human-AI systems (Ferrario et al., 2026): rubric satisfaction can function as an interpretable reliability indicator for tutoring processes.

### Limitations  
- Our teacher feedback and rubrics, while grounded in ConvoLearn dimensions, may not capture all classroom realities (e.g., affect, accessibility accommodations).  
- Process reward modeling depends on critique quality; scaling human critiques remains challenging even with MetaRM-style generalization (Wang et al., 2026).  
- We did not test cross-domain transfer beyond Earth Science; future work should evaluate other subjects and age groups.

---

## Conclusion  
We presented **PACE-Tutor**, a hybrid framework for constructivist LLM tutoring that extends ConvoLearn-style instruction tuning with (i) **agentic context evolution** to reduce retrieval noise and cost, and (ii) **process-aligned reinforcement learning** using natural-language teacher feedback, stabilized by **segmental advantage estimation**. Across rubric-based constructivist evaluation, PACE-Tutor improves pedagogical quality and reduces solution leakage while maintaining efficient context use. The work suggests that scaling constructivist tutoring requires not only better datasets, but also process-grounded rewards, adaptive context strategies, and fine-grained evaluation.

---

## Contributions (explicit)  
1. **PACE-Tutor framework:** A combined instruction-tuning + ACE-based context evolution + process-aligned RL pipeline for constructivist tutoring grounded in ConvoLearn (Sharma et al., 2026; Chen et al., 2026; Wang et al., 2026).  
2. **Process reward adaptation for tutoring:** A tutoring-specific instantiation of RM-NLHF-style natural-language feedback rewards and MetaRM-style generalization (Wang et al., 2026).  
3. **Long-dialogue RL stabilization:** Application of SAE to multi-turn tutoring RL to improve stability and outcomes (Gong et al., 2026).  
4. **Constructivist Rubric Set (CRS):** A fine-grained, verifiable rubric evaluation protocol inspired by DeepResearch Bench II, aligned to ConvoLearn’s six pedagogical dimensions (Li et al., 2026; Sharma et al., 2026).  
5. **Empirical ablations:** Evidence that (a) ACE primarily improves efficiency and reduces irrelevant retrieval, while (b) process-aligned RL drives the largest gains in constructivist rubric satisfaction and leakage reduction.

---